% ============================================================
%  35_Model.tex  --  Model opisowy i matematyczny wyznaczania
%                    danych referencyjnych (wersja uproszczona)
%  Wstawiany do rozdzialu 3 poleceniem:  % ============================================================
%  35_Model.tex  --  Model opisowy i matematyczny wyznaczania
%                    danych referencyjnych (wersja uproszczona)
%  Wstawiany do rozdzialu 3 poleceniem:  % ============================================================
%  35_Model.tex  --  Model opisowy i matematyczny wyznaczania
%                    danych referencyjnych (wersja uproszczona)
%  Wstawiany do rozdzialu 3 poleceniem:  % ============================================================
%  35_Model.tex  --  Model opisowy i matematyczny wyznaczania
%                    danych referencyjnych (wersja uproszczona)
%  Wstawiany do rozdzialu 3 poleceniem:  \input{tex/35_Model}
%  Wymaga w preambule: amsmath, siunitx, enumitem
% ============================================================

% Objaśnienia symboli pod wzorem ("gdzie:"): symbol w kolumnie od
% wcięcia akapitu, myślniki zawsze w jednej pionowej linii
\newenvironment{objasnienia}{%
  \noindent gdzie:
  \begin{description}[font=\normalfont, labelindent=\parindent,
    labelwidth=1.2cm, labelsep=0.2cm, leftmargin=!, align=left,
    itemsep=0pt, topsep=2pt]}{\end{description}}

\section{Model wyznaczania danych referencyjnych}
\label{sec:model}

Zarejestrowane przez stację dane mają posłużyć jako dane referencyjne, czyli
mają pozwolić stwierdzić, gdzie znajdował się dany statek powietrzny
w określonej chwili czasu. % referencja do podrozdziału.
W tym celu potrzebny jest model, który łączy trzy źródła informacji:
komunikaty ADS-B odbierane drogą radiową, pozycję stacji z odbiornika GPS oraz
zegar systemowy komputera, którym oznaczane są wszystkie zapisane wiadomości.
W niniejszym podrozdziale model ten przedstawiono najpierw w postaci słownej
(model opisowy), a następnie w postaci wzorów (model matematyczny). Na końcu
omówiono przyjęte uproszczenia i podano przykład obliczeniowy.

\subsection{Model opisowy}
\label{sec:model_opisowy}

Statek powietrzny wyposażony w transponder ADS-B kilka razy na sekundę
rozgłasza na częstotliwości \SI{1090}{MHz} komunikaty zawierające m.in.\ jego
unikalny adres ICAO, pozycję (szerokość i długość geograficzną) wyznaczoną
przez pokładowy odbiornik GNSS, wysokość barometryczną oraz prędkość
względem ziemi~\cite{sun2021}. Stacja odbiera część tych komunikatów --
niektóre giną z powodu zakłóceń, zbyt słabego sygnału lub przesłonięcia
anteny. Odebrane komunikaty są zapisywane w postaci tekstowej (format SBS)
razem ze znacznikiem czasu nadawanym przez komputer w chwili odbioru.

Komunikat ADS-B mówi, gdzie samolot jest, ale nie mówi, gdzie jest
względem stacji. Do tego potrzebna jest pozycja samej stacji, którą
podaje odbiornik GPS w komunikacie NMEA \texttt{GGA}~\cite{nmea0183}.
Znając obie pozycje, można obliczyć odległość od stacji do samolotu.

Drugim problemem jest czas. Komputer Raspberry Pi nie ma podtrzymywanego
bateryjnie zegara, a w terenie nie ma dostępu do internetu, więc jego zegar
może się mylić o nieznaną wartość. Ponieważ wszystkie zapisane wiadomości są
oznaczone właśnie tym zegarem, błąd zegara przenosi się na wszystkie dane.
Odbiornik GPS co sekundę przesyła jednak dokładny czas UTC (komunikat
\texttt{RMC}). Komputer zapisuje go razem z własnym wskazaniem zegara, dzięki
czemu po zakończeniu pomiarów można zmierzyć, o ile zegar się mylił,
i poprawić czas wszystkich wiadomości.

Granice modelowanego układu wyznacza sama stacja (antena, odbiornik RTL-SDR,
odbiornik GPS i komputer). Samoloty i satelity GPS należą do otoczenia,
a punktami styku są: łącze radiowe \SI{1090}{MHz}, łącze szeregowe
odbiornika GPS oraz plik zapisu rosbag2. Model składa się z dwóch części:
\begin{enumerate}[nosep]
  \item modelu czasu -- poprawienie znaczników czasu zegara systemowego
        na czas UTC (podrozdział~\ref{sec:model_czasu}),
  \item modelu odległości -- obliczenie odległości od stacji do samolotu
        (podrozdział~\ref{sec:model_odleglosc}).
\end{enumerate}

\subsection{Model czasu}
\label{sec:model_czasu}

Niech $t$ oznacza wskazanie zegara systemowego komputera, a $\tau$ --
odpowiadający mu rzeczywisty czas UTC. Zależność między nimi zapisano jako
\begin{equation}
  \tau = t + \delta,
  \label{eq:model_zegara}
\end{equation}
gdzie $\delta$ jest przesunięciem (błędem) zegara systemowego w sekundach.
Przyjęto, że w trakcie jednej sesji pomiarowej przesunięcie $\delta$ jest
stałe.

Węzeł obsługi GPS publikuje w temacie \texttt{/gps/time\_ref} pary
$(t_k, \tau_k)$, $k = 1, 2, \ldots, K$, gdzie $\tau_k$ jest czasem UTC
odczytanym z komunikatu \texttt{RMC}, a $t_k$ -- wskazaniem zegara
systemowego w chwili odczytu tego komunikatu. Każda para daje jeden pomiar
przesunięcia:
\begin{equation}
  \delta_k = \tau_k - t_k, \qquad k = 1, 2, \ldots, K.
  \label{eq:offset_pomiar}
\end{equation}
Pojedyncze pomiary różnią się nieco od siebie, dlatego jako wynik przyjmuje
się ich średnią arytmetyczną, a rozrzut pomiarów opisuje się odchyleniem
standardowym:
\begin{equation}
  \hat{\delta} = \frac{1}{K} \sum_{k=1}^{K} \delta_k,
  \qquad
  \sigma_{\delta} = \sqrt{\frac{1}{K-1} \sum_{k=1}^{K}
                    \bigl(\delta_k - \hat{\delta}\bigr)^2},
  \label{eq:offset_srednia}
\end{equation}
\begin{objasnienia}
  \item[$\hat{\delta}$] -- wyznaczone przesunięcie zegara w sekundach,
  \item[$\sigma_{\delta}$] -- odchylenie standardowe pomiarów przesunięcia w sekundach,
  \item[$K$] -- liczba pomiarów (przy jednym pomiarze na sekundę sesja trwająca \SI{10}{min} daje ok.~600 pomiarów).
\end{objasnienia}
Poprawiony czas UTC dowolnej zapisanej wiadomości o znaczniku $t$ wynosi
\begin{equation}
  \hat{\tau} = t + \hat{\delta}.
  \label{eq:korekta_czasu}
\end{equation}

Pomiar \eqref{eq:offset_pomiar} nie jest idealny, ponieważ znacznik $t_k$
jest nadawany dopiero po przesłaniu komunikatu łączem szeregowym i odczytaniu
go przez program. Przy prędkości \num{9600}~bit/s przesłanie jednego znaku
(10~bitów) trwa ok.\ \SI{1.04}{ms}, a komunikat NMEA ma do 82 znaków, więc
sama transmisja zajmuje do ok.\ \SI{85}{ms}. Program odczytuje port co
\SI{100}{ms}, co dodaje średnio kolejne \SI{50}{ms}. Wyznaczone przesunięcie
$\hat{\delta}$ jest więc obarczone błędem rzędu \SI{0.1}{s}, co -- jak
pokazano w podrozdziale~\ref{sec:model_przyklad} -- jest wartością
akceptowalną.

\subsection{Model odległości}
\label{sec:model_odleglosc}

\paragraph{Dane wejściowe.}
Z komunikatu SBS typu \texttt{MSG,3} odczytuje się szerokość $\varphi_i$
i długość $\lambda_i$ geograficzną samolotu (w stopniach) oraz jego wysokość
w stopach. Z tematu \texttt{/gps/fix} odczytuje się szerokość
$\varphi_{\mathrm{s}}$, długość $\lambda_{\mathrm{s}}$ i wysokość
$h_{\mathrm{s}}$ stacji (w metrach). Indeks $i$ oznacza samolot, a indeks
$\mathrm{s}$ -- stację.

\paragraph{Przeliczenie jednostek.}
W lotnictwie wysokość podaje się w stopach, a prędkość w węzłach. Na
jednostki układu SI przelicza się je wzorami
\begin{equation}
  h_i = \num{0.3048}\, h_{\mathrm{ft}}, \qquad
  v = \num{0.514444}\, v_{\mathrm{kt}},
  \label{eq:jednostki}
\end{equation}
\begin{objasnienia}
  \item[$h_{\mathrm{ft}}$] -- wysokość samolotu w stopach,
  \item[$h_i$] -- wysokość samolotu w metrach,
  \item[$v_{\mathrm{kt}}$] -- prędkość w węzłach,
  \item[$v$] -- prędkość w metrach na sekundę.
\end{objasnienia}

\paragraph{Odległość po powierzchni Ziemi.}
Ziemię przybliżono kulą o promieniu $R_{\mathrm{Z}} = \SI{6371}{km}$.
Odległość między stacją a punktem na powierzchni Ziemi leżącym dokładnie pod
samolotem oblicza się wzorem haversine:
\begin{equation}
  s = 2 R_{\mathrm{Z}} \arcsin \sqrt{
        \sin^2\frac{\Delta\varphi}{2}
        + \cos\varphi_{\mathrm{s}} \cos\varphi_i
          \sin^2\frac{\Delta\lambda}{2}},
  \label{eq:haversine}
\end{equation}
\begin{objasnienia}
  \item[$s$] -- odległość mierzona po powierzchni Ziemi w metrach,
  \item[$\Delta\varphi$] -- różnica szerokości geograficznych, $\Delta\varphi = \varphi_i - \varphi_{\mathrm{s}}$, w radianach,
  \item[$\Delta\lambda$] -- różnica długości geograficznych, $\Delta\lambda = \lambda_i - \lambda_{\mathrm{s}}$, w radianach,
  \item[$R_{\mathrm{Z}}$] -- średni promień Ziemi w metrach.
\end{objasnienia}
Kąty podawane w stopniach należy przed użyciem we wzorze
\eqref{eq:haversine} zamienić na radiany, mnożąc je przez $\pi/180$.

\paragraph{Odległość skośna.}
Samolot znajduje się na wysokości $h_i$, a stacja na wysokości
$h_{\mathrm{s}}$. Odległość w linii prostej od stacji do samolotu
(odległość skośną) przybliża się, traktując odległość $s$ i różnicę wysokości
jako przyprostokątne trójkąta prostokątnego:
\begin{equation}
  \rho = \sqrt{s^2 + \Delta h^2}, \qquad \Delta h = h_i - h_{\mathrm{s}},
  \label{eq:odleglosc_skosna}
\end{equation}
gdzie $\rho$ jest odległością skośną, a $\Delta h$ -- różnicą wysokości
samolotu i stacji, obie w metrach. Wzory \eqref{eq:jednostki}%
--\eqref{eq:odleglosc_skosna} są zaimplementowane w węźle wyświetlacza
stacji, który na ich podstawie pokazuje odległość do każdego samolotu.

\subsection{Uproszczenia modelu}
\label{sec:model_uproszczenia}

W modelu przyjęto następujące uproszczenia:
\begin{itemize}[nosep]
  \item Ziemia jest kulą, a nie elipsoidą -- powoduje to błąd odległości
        rzędu kilku dziesiątych procenta,
  \item we wzorze \eqref{eq:odleglosc_skosna} pominięto krzywiznę Ziemi;
        błąd rośnie wraz z odległością samolotu od stacji,
  \item wysokość barometryczna z ADS-B jest traktowana jak wysokość nad
        poziomem morza; przy ciśnieniu innym niż standardowe
        (\SI{1013.25}{hPa}) różnica może sięgać kilkudziesięciu metrów,
  \item przesunięcie zegara $\delta$ jest stałe w trakcie sesji (pominięto
        powolny dryf zegara),
  \item stacja jest nieruchoma w trakcie sesji,
  \item pominięto opóźnienie między odebraniem komunikatu ADS-B a nadaniem mu
        znacznika czasu (pojedyncze milisekundy).
\end{itemize}
Uproszczenia te są wystarczające do podglądu na żywo i do analizy zasięgu
stacji. Przy dalszym rozwoju systemu model odległości można zastąpić
dokładniejszym przeliczeniem współrzędnych do lokalnego układu kartezjańskiego
związanego ze stacją, które uwzględnia kształt Ziemi i pozwala dodatkowo
wyznaczyć kierunek obserwacji samolotu.

\subsection{Przykład obliczeniowy}
\label{sec:model_przyklad}

Przyjęto stację w punkcie $\varphi_{\mathrm{s}} = \SI{50.0}{\degree}$,
$\lambda_{\mathrm{s}} = \SI{20.0}{\degree}$, $h_{\mathrm{s}} = \SI{250}{m}$
oraz samolot w punkcie $\varphi_i = \SI{50.3169}{\degree}$,
$\lambda_i = \SI{20.49795}{\degree}$, lecący na wysokości
$h_{\mathrm{ft}} = \num{32800}$~stóp.

Z wzoru \eqref{eq:jednostki} wysokość samolotu wynosi
\begin{equation}
  h_i = \num{0.3048} \cdot \num{32800} \approx \SI{9997.4}{m},
  \qquad
  \Delta h = \num{9997.4} - 250 = \SI{9747.4}{m}.
\end{equation}
Różnice współrzędnych wynoszą $\Delta\varphi = \SI{0.3169}{\degree}$
i $\Delta\lambda = \SI{0.49795}{\degree}$. Wyrażenie pod pierwiastkiem we
wzorze \eqref{eq:haversine} ma wartość ok.\ \num{1.540e-5}, skąd
\begin{equation}
  s \approx \SI{50.00}{km}, \qquad
  \rho = \sqrt{\num{50000}^2 + \num{9747.4}^2} \approx \SI{50.94}{km}.
\end{equation}

Wpływ błędu czasu pokazuje następujące oszacowanie: samolot lecący
z prędkością $v = \SI{230}{m/s}$ w czasie $\Delta\tau$ pokonuje drogę
\begin{equation}
  \Delta x = v \, \Delta\tau .
  \label{eq:blad_czasu}
\end{equation}
Jeżeli zegar stacji myli się o \SI{10}{s} i nie zostanie poprawiony wzorem
\eqref{eq:korekta_czasu}, pozycja przypisana danej chwili będzie błędna
o ok.\ \SI{2.3}{km}. Po korekcji błąd czasu spada do ok.\ \SI{0.1}{s},
a błąd pozycji -- do ok.\ \SI{23}{m}. Uzasadnia to rejestrowanie czasu GPS
razem z danymi ADS-B.

\subsection{Weryfikacja modelu}
\label{sec:model_weryfikacja}

Poprawność modelu odległości można sprawdzić na prostych przypadkach:
dla samolotu położonego dokładnie nad stacją powinno zachodzić $s = 0$
i $\rho = \Delta h$, a dla punktów odległych o \SI{1}{\degree} szerokości
geograficznej $s \approx \SI{111.2}{km}$. Poprawność modelu czasu można
sprawdzić, porównując $\hat{\delta}$ wyznaczone w sesji z dostępem do
internetu (gdy zegar jest synchronizowany przez NTP i $\hat{\delta}$ powinno
być bliskie zera) z sesją bez dostępu do sieci. Szczegółową ocenę na
podstawie zarejestrowanych danych przedstawiono w~rozdziale~\ref{chap:badania}.
%  Wymaga w preambule: amsmath, siunitx, enumitem
% ============================================================

% Objaśnienia symboli pod wzorem ("gdzie:"): symbol w kolumnie od
% wcięcia akapitu, myślniki zawsze w jednej pionowej linii
\newenvironment{objasnienia}{%
  \noindent gdzie:
  \begin{description}[font=\normalfont, labelindent=\parindent,
    labelwidth=1.2cm, labelsep=0.2cm, leftmargin=!, align=left,
    itemsep=0pt, topsep=2pt]}{\end{description}}

\section{Model wyznaczania danych referencyjnych}
\label{sec:model}

Zarejestrowane przez stację dane mają posłużyć jako dane referencyjne, czyli
mają pozwolić stwierdzić, gdzie znajdował się dany statek powietrzny
w określonej chwili czasu. % referencja do podrozdziału.
W tym celu potrzebny jest model, który łączy trzy źródła informacji:
komunikaty ADS-B odbierane drogą radiową, pozycję stacji z odbiornika GPS oraz
zegar systemowy komputera, którym oznaczane są wszystkie zapisane wiadomości.
W niniejszym podrozdziale model ten przedstawiono najpierw w postaci słownej
(model opisowy), a następnie w postaci wzorów (model matematyczny). Na końcu
omówiono przyjęte uproszczenia i podano przykład obliczeniowy.

\subsection{Model opisowy}
\label{sec:model_opisowy}

Statek powietrzny wyposażony w transponder ADS-B kilka razy na sekundę
rozgłasza na częstotliwości \SI{1090}{MHz} komunikaty zawierające m.in.\ jego
unikalny adres ICAO, pozycję (szerokość i długość geograficzną) wyznaczoną
przez pokładowy odbiornik GNSS, wysokość barometryczną oraz prędkość
względem ziemi~\cite{sun2021}. Stacja odbiera część tych komunikatów --
niektóre giną z powodu zakłóceń, zbyt słabego sygnału lub przesłonięcia
anteny. Odebrane komunikaty są zapisywane w postaci tekstowej (format SBS)
razem ze znacznikiem czasu nadawanym przez komputer w chwili odbioru.

Komunikat ADS-B mówi, gdzie samolot jest, ale nie mówi, gdzie jest
względem stacji. Do tego potrzebna jest pozycja samej stacji, którą
podaje odbiornik GPS w komunikacie NMEA \texttt{GGA}~\cite{nmea0183}.
Znając obie pozycje, można obliczyć odległość od stacji do samolotu.

Drugim problemem jest czas. Komputer Raspberry Pi nie ma podtrzymywanego
bateryjnie zegara, a w terenie nie ma dostępu do internetu, więc jego zegar
może się mylić o nieznaną wartość. Ponieważ wszystkie zapisane wiadomości są
oznaczone właśnie tym zegarem, błąd zegara przenosi się na wszystkie dane.
Odbiornik GPS co sekundę przesyła jednak dokładny czas UTC (komunikat
\texttt{RMC}). Komputer zapisuje go razem z własnym wskazaniem zegara, dzięki
czemu po zakończeniu pomiarów można zmierzyć, o ile zegar się mylił,
i poprawić czas wszystkich wiadomości.

Granice modelowanego układu wyznacza sama stacja (antena, odbiornik RTL-SDR,
odbiornik GPS i komputer). Samoloty i satelity GPS należą do otoczenia,
a punktami styku są: łącze radiowe \SI{1090}{MHz}, łącze szeregowe
odbiornika GPS oraz plik zapisu rosbag2. Model składa się z dwóch części:
\begin{enumerate}[nosep]
  \item modelu czasu -- poprawienie znaczników czasu zegara systemowego
        na czas UTC (podrozdział~\ref{sec:model_czasu}),
  \item modelu odległości -- obliczenie odległości od stacji do samolotu
        (podrozdział~\ref{sec:model_odleglosc}).
\end{enumerate}

\subsection{Model czasu}
\label{sec:model_czasu}

Niech $t$ oznacza wskazanie zegara systemowego komputera, a $\tau$ --
odpowiadający mu rzeczywisty czas UTC. Zależność między nimi zapisano jako
\begin{equation}
  \tau = t + \delta,
  \label{eq:model_zegara}
\end{equation}
gdzie $\delta$ jest przesunięciem (błędem) zegara systemowego w sekundach.
Przyjęto, że w trakcie jednej sesji pomiarowej przesunięcie $\delta$ jest
stałe.

Węzeł obsługi GPS publikuje w temacie \texttt{/gps/time\_ref} pary
$(t_k, \tau_k)$, $k = 1, 2, \ldots, K$, gdzie $\tau_k$ jest czasem UTC
odczytanym z komunikatu \texttt{RMC}, a $t_k$ -- wskazaniem zegara
systemowego w chwili odczytu tego komunikatu. Każda para daje jeden pomiar
przesunięcia:
\begin{equation}
  \delta_k = \tau_k - t_k, \qquad k = 1, 2, \ldots, K.
  \label{eq:offset_pomiar}
\end{equation}
Pojedyncze pomiary różnią się nieco od siebie, dlatego jako wynik przyjmuje
się ich średnią arytmetyczną, a rozrzut pomiarów opisuje się odchyleniem
standardowym:
\begin{equation}
  \hat{\delta} = \frac{1}{K} \sum_{k=1}^{K} \delta_k,
  \qquad
  \sigma_{\delta} = \sqrt{\frac{1}{K-1} \sum_{k=1}^{K}
                    \bigl(\delta_k - \hat{\delta}\bigr)^2},
  \label{eq:offset_srednia}
\end{equation}
\begin{objasnienia}
  \item[$\hat{\delta}$] -- wyznaczone przesunięcie zegara w sekundach,
  \item[$\sigma_{\delta}$] -- odchylenie standardowe pomiarów przesunięcia w sekundach,
  \item[$K$] -- liczba pomiarów (przy jednym pomiarze na sekundę sesja trwająca \SI{10}{min} daje ok.~600 pomiarów).
\end{objasnienia}
Poprawiony czas UTC dowolnej zapisanej wiadomości o znaczniku $t$ wynosi
\begin{equation}
  \hat{\tau} = t + \hat{\delta}.
  \label{eq:korekta_czasu}
\end{equation}

Pomiar \eqref{eq:offset_pomiar} nie jest idealny, ponieważ znacznik $t_k$
jest nadawany dopiero po przesłaniu komunikatu łączem szeregowym i odczytaniu
go przez program. Przy prędkości \num{9600}~bit/s przesłanie jednego znaku
(10~bitów) trwa ok.\ \SI{1.04}{ms}, a komunikat NMEA ma do 82 znaków, więc
sama transmisja zajmuje do ok.\ \SI{85}{ms}. Program odczytuje port co
\SI{100}{ms}, co dodaje średnio kolejne \SI{50}{ms}. Wyznaczone przesunięcie
$\hat{\delta}$ jest więc obarczone błędem rzędu \SI{0.1}{s}, co -- jak
pokazano w podrozdziale~\ref{sec:model_przyklad} -- jest wartością
akceptowalną.

\subsection{Model odległości}
\label{sec:model_odleglosc}

\paragraph{Dane wejściowe.}
Z komunikatu SBS typu \texttt{MSG,3} odczytuje się szerokość $\varphi_i$
i długość $\lambda_i$ geograficzną samolotu (w stopniach) oraz jego wysokość
w stopach. Z tematu \texttt{/gps/fix} odczytuje się szerokość
$\varphi_{\mathrm{s}}$, długość $\lambda_{\mathrm{s}}$ i wysokość
$h_{\mathrm{s}}$ stacji (w metrach). Indeks $i$ oznacza samolot, a indeks
$\mathrm{s}$ -- stację.

\paragraph{Przeliczenie jednostek.}
W lotnictwie wysokość podaje się w stopach, a prędkość w węzłach. Na
jednostki układu SI przelicza się je wzorami
\begin{equation}
  h_i = \num{0.3048}\, h_{\mathrm{ft}}, \qquad
  v = \num{0.514444}\, v_{\mathrm{kt}},
  \label{eq:jednostki}
\end{equation}
\begin{objasnienia}
  \item[$h_{\mathrm{ft}}$] -- wysokość samolotu w stopach,
  \item[$h_i$] -- wysokość samolotu w metrach,
  \item[$v_{\mathrm{kt}}$] -- prędkość w węzłach,
  \item[$v$] -- prędkość w metrach na sekundę.
\end{objasnienia}

\paragraph{Odległość po powierzchni Ziemi.}
Ziemię przybliżono kulą o promieniu $R_{\mathrm{Z}} = \SI{6371}{km}$.
Odległość między stacją a punktem na powierzchni Ziemi leżącym dokładnie pod
samolotem oblicza się wzorem haversine:
\begin{equation}
  s = 2 R_{\mathrm{Z}} \arcsin \sqrt{
        \sin^2\frac{\Delta\varphi}{2}
        + \cos\varphi_{\mathrm{s}} \cos\varphi_i
          \sin^2\frac{\Delta\lambda}{2}},
  \label{eq:haversine}
\end{equation}
\begin{objasnienia}
  \item[$s$] -- odległość mierzona po powierzchni Ziemi w metrach,
  \item[$\Delta\varphi$] -- różnica szerokości geograficznych, $\Delta\varphi = \varphi_i - \varphi_{\mathrm{s}}$, w radianach,
  \item[$\Delta\lambda$] -- różnica długości geograficznych, $\Delta\lambda = \lambda_i - \lambda_{\mathrm{s}}$, w radianach,
  \item[$R_{\mathrm{Z}}$] -- średni promień Ziemi w metrach.
\end{objasnienia}
Kąty podawane w stopniach należy przed użyciem we wzorze
\eqref{eq:haversine} zamienić na radiany, mnożąc je przez $\pi/180$.

\paragraph{Odległość skośna.}
Samolot znajduje się na wysokości $h_i$, a stacja na wysokości
$h_{\mathrm{s}}$. Odległość w linii prostej od stacji do samolotu
(odległość skośną) przybliża się, traktując odległość $s$ i różnicę wysokości
jako przyprostokątne trójkąta prostokątnego:
\begin{equation}
  \rho = \sqrt{s^2 + \Delta h^2}, \qquad \Delta h = h_i - h_{\mathrm{s}},
  \label{eq:odleglosc_skosna}
\end{equation}
gdzie $\rho$ jest odległością skośną, a $\Delta h$ -- różnicą wysokości
samolotu i stacji, obie w metrach. Wzory \eqref{eq:jednostki}%
--\eqref{eq:odleglosc_skosna} są zaimplementowane w węźle wyświetlacza
stacji, który na ich podstawie pokazuje odległość do każdego samolotu.

\subsection{Uproszczenia modelu}
\label{sec:model_uproszczenia}

W modelu przyjęto następujące uproszczenia:
\begin{itemize}[nosep]
  \item Ziemia jest kulą, a nie elipsoidą -- powoduje to błąd odległości
        rzędu kilku dziesiątych procenta,
  \item we wzorze \eqref{eq:odleglosc_skosna} pominięto krzywiznę Ziemi;
        błąd rośnie wraz z odległością samolotu od stacji,
  \item wysokość barometryczna z ADS-B jest traktowana jak wysokość nad
        poziomem morza; przy ciśnieniu innym niż standardowe
        (\SI{1013.25}{hPa}) różnica może sięgać kilkudziesięciu metrów,
  \item przesunięcie zegara $\delta$ jest stałe w trakcie sesji (pominięto
        powolny dryf zegara),
  \item stacja jest nieruchoma w trakcie sesji,
  \item pominięto opóźnienie między odebraniem komunikatu ADS-B a nadaniem mu
        znacznika czasu (pojedyncze milisekundy).
\end{itemize}
Uproszczenia te są wystarczające do podglądu na żywo i do analizy zasięgu
stacji. Przy dalszym rozwoju systemu model odległości można zastąpić
dokładniejszym przeliczeniem współrzędnych do lokalnego układu kartezjańskiego
związanego ze stacją, które uwzględnia kształt Ziemi i pozwala dodatkowo
wyznaczyć kierunek obserwacji samolotu.

\subsection{Przykład obliczeniowy}
\label{sec:model_przyklad}

Przyjęto stację w punkcie $\varphi_{\mathrm{s}} = \SI{50.0}{\degree}$,
$\lambda_{\mathrm{s}} = \SI{20.0}{\degree}$, $h_{\mathrm{s}} = \SI{250}{m}$
oraz samolot w punkcie $\varphi_i = \SI{50.3169}{\degree}$,
$\lambda_i = \SI{20.49795}{\degree}$, lecący na wysokości
$h_{\mathrm{ft}} = \num{32800}$~stóp.

Z wzoru \eqref{eq:jednostki} wysokość samolotu wynosi
\begin{equation}
  h_i = \num{0.3048} \cdot \num{32800} \approx \SI{9997.4}{m},
  \qquad
  \Delta h = \num{9997.4} - 250 = \SI{9747.4}{m}.
\end{equation}
Różnice współrzędnych wynoszą $\Delta\varphi = \SI{0.3169}{\degree}$
i $\Delta\lambda = \SI{0.49795}{\degree}$. Wyrażenie pod pierwiastkiem we
wzorze \eqref{eq:haversine} ma wartość ok.\ \num{1.540e-5}, skąd
\begin{equation}
  s \approx \SI{50.00}{km}, \qquad
  \rho = \sqrt{\num{50000}^2 + \num{9747.4}^2} \approx \SI{50.94}{km}.
\end{equation}

Wpływ błędu czasu pokazuje następujące oszacowanie: samolot lecący
z prędkością $v = \SI{230}{m/s}$ w czasie $\Delta\tau$ pokonuje drogę
\begin{equation}
  \Delta x = v \, \Delta\tau .
  \label{eq:blad_czasu}
\end{equation}
Jeżeli zegar stacji myli się o \SI{10}{s} i nie zostanie poprawiony wzorem
\eqref{eq:korekta_czasu}, pozycja przypisana danej chwili będzie błędna
o ok.\ \SI{2.3}{km}. Po korekcji błąd czasu spada do ok.\ \SI{0.1}{s},
a błąd pozycji -- do ok.\ \SI{23}{m}. Uzasadnia to rejestrowanie czasu GPS
razem z danymi ADS-B.

\subsection{Weryfikacja modelu}
\label{sec:model_weryfikacja}

Poprawność modelu odległości można sprawdzić na prostych przypadkach:
dla samolotu położonego dokładnie nad stacją powinno zachodzić $s = 0$
i $\rho = \Delta h$, a dla punktów odległych o \SI{1}{\degree} szerokości
geograficznej $s \approx \SI{111.2}{km}$. Poprawność modelu czasu można
sprawdzić, porównując $\hat{\delta}$ wyznaczone w sesji z dostępem do
internetu (gdy zegar jest synchronizowany przez NTP i $\hat{\delta}$ powinno
być bliskie zera) z sesją bez dostępu do sieci. Szczegółową ocenę na
podstawie zarejestrowanych danych przedstawiono w~rozdziale~\ref{chap:badania}.
%  Wymaga w preambule: amsmath, siunitx, enumitem
% ============================================================

% Objaśnienia symboli pod wzorem ("gdzie:"): symbol w kolumnie od
% wcięcia akapitu, myślniki zawsze w jednej pionowej linii
\newenvironment{objasnienia}{%
  \noindent gdzie:
  \begin{description}[font=\normalfont, labelindent=\parindent,
    labelwidth=1.2cm, labelsep=0.2cm, leftmargin=!, align=left,
    itemsep=0pt, topsep=2pt]}{\end{description}}

\section{Model wyznaczania danych referencyjnych}
\label{sec:model}

Zarejestrowane przez stację dane mają posłużyć jako dane referencyjne, czyli
mają pozwolić stwierdzić, gdzie znajdował się dany statek powietrzny
w określonej chwili czasu. % referencja do podrozdziału.
W tym celu potrzebny jest model, który łączy trzy źródła informacji:
komunikaty ADS-B odbierane drogą radiową, pozycję stacji z odbiornika GPS oraz
zegar systemowy komputera, którym oznaczane są wszystkie zapisane wiadomości.
W niniejszym podrozdziale model ten przedstawiono najpierw w postaci słownej
(model opisowy), a następnie w postaci wzorów (model matematyczny). Na końcu
omówiono przyjęte uproszczenia i podano przykład obliczeniowy.

\subsection{Model opisowy}
\label{sec:model_opisowy}

Statek powietrzny wyposażony w transponder ADS-B kilka razy na sekundę
rozgłasza na częstotliwości \SI{1090}{MHz} komunikaty zawierające m.in.\ jego
unikalny adres ICAO, pozycję (szerokość i długość geograficzną) wyznaczoną
przez pokładowy odbiornik GNSS, wysokość barometryczną oraz prędkość
względem ziemi~\cite{sun2021}. Stacja odbiera część tych komunikatów --
niektóre giną z powodu zakłóceń, zbyt słabego sygnału lub przesłonięcia
anteny. Odebrane komunikaty są zapisywane w postaci tekstowej (format SBS)
razem ze znacznikiem czasu nadawanym przez komputer w chwili odbioru.

Komunikat ADS-B mówi, gdzie samolot jest, ale nie mówi, gdzie jest
względem stacji. Do tego potrzebna jest pozycja samej stacji, którą
podaje odbiornik GPS w komunikacie NMEA \texttt{GGA}~\cite{nmea0183}.
Znając obie pozycje, można obliczyć odległość od stacji do samolotu.

Drugim problemem jest czas. Komputer Raspberry Pi nie ma podtrzymywanego
bateryjnie zegara, a w terenie nie ma dostępu do internetu, więc jego zegar
może się mylić o nieznaną wartość. Ponieważ wszystkie zapisane wiadomości są
oznaczone właśnie tym zegarem, błąd zegara przenosi się na wszystkie dane.
Odbiornik GPS co sekundę przesyła jednak dokładny czas UTC (komunikat
\texttt{RMC}). Komputer zapisuje go razem z własnym wskazaniem zegara, dzięki
czemu po zakończeniu pomiarów można zmierzyć, o ile zegar się mylił,
i poprawić czas wszystkich wiadomości.

Granice modelowanego układu wyznacza sama stacja (antena, odbiornik RTL-SDR,
odbiornik GPS i komputer). Samoloty i satelity GPS należą do otoczenia,
a punktami styku są: łącze radiowe \SI{1090}{MHz}, łącze szeregowe
odbiornika GPS oraz plik zapisu rosbag2. Model składa się z dwóch części:
\begin{enumerate}[nosep]
  \item modelu czasu -- poprawienie znaczników czasu zegara systemowego
        na czas UTC (podrozdział~\ref{sec:model_czasu}),
  \item modelu odległości -- obliczenie odległości od stacji do samolotu
        (podrozdział~\ref{sec:model_odleglosc}).
\end{enumerate}

\subsection{Model czasu}
\label{sec:model_czasu}

Niech $t$ oznacza wskazanie zegara systemowego komputera, a $\tau$ --
odpowiadający mu rzeczywisty czas UTC. Zależność między nimi zapisano jako
\begin{equation}
  \tau = t + \delta,
  \label{eq:model_zegara}
\end{equation}
gdzie $\delta$ jest przesunięciem (błędem) zegara systemowego w sekundach.
Przyjęto, że w trakcie jednej sesji pomiarowej przesunięcie $\delta$ jest
stałe.

Węzeł obsługi GPS publikuje w temacie \texttt{/gps/time\_ref} pary
$(t_k, \tau_k)$, $k = 1, 2, \ldots, K$, gdzie $\tau_k$ jest czasem UTC
odczytanym z komunikatu \texttt{RMC}, a $t_k$ -- wskazaniem zegara
systemowego w chwili odczytu tego komunikatu. Każda para daje jeden pomiar
przesunięcia:
\begin{equation}
  \delta_k = \tau_k - t_k, \qquad k = 1, 2, \ldots, K.
  \label{eq:offset_pomiar}
\end{equation}
Pojedyncze pomiary różnią się nieco od siebie, dlatego jako wynik przyjmuje
się ich średnią arytmetyczną, a rozrzut pomiarów opisuje się odchyleniem
standardowym:
\begin{equation}
  \hat{\delta} = \frac{1}{K} \sum_{k=1}^{K} \delta_k,
  \qquad
  \sigma_{\delta} = \sqrt{\frac{1}{K-1} \sum_{k=1}^{K}
                    \bigl(\delta_k - \hat{\delta}\bigr)^2},
  \label{eq:offset_srednia}
\end{equation}
\begin{objasnienia}
  \item[$\hat{\delta}$] -- wyznaczone przesunięcie zegara w sekundach,
  \item[$\sigma_{\delta}$] -- odchylenie standardowe pomiarów przesunięcia w sekundach,
  \item[$K$] -- liczba pomiarów (przy jednym pomiarze na sekundę sesja trwająca \SI{10}{min} daje ok.~600 pomiarów).
\end{objasnienia}
Poprawiony czas UTC dowolnej zapisanej wiadomości o znaczniku $t$ wynosi
\begin{equation}
  \hat{\tau} = t + \hat{\delta}.
  \label{eq:korekta_czasu}
\end{equation}

Pomiar \eqref{eq:offset_pomiar} nie jest idealny, ponieważ znacznik $t_k$
jest nadawany dopiero po przesłaniu komunikatu łączem szeregowym i odczytaniu
go przez program. Przy prędkości \num{9600}~bit/s przesłanie jednego znaku
(10~bitów) trwa ok.\ \SI{1.04}{ms}, a komunikat NMEA ma do 82 znaków, więc
sama transmisja zajmuje do ok.\ \SI{85}{ms}. Program odczytuje port co
\SI{100}{ms}, co dodaje średnio kolejne \SI{50}{ms}. Wyznaczone przesunięcie
$\hat{\delta}$ jest więc obarczone błędem rzędu \SI{0.1}{s}, co -- jak
pokazano w podrozdziale~\ref{sec:model_przyklad} -- jest wartością
akceptowalną.

\subsection{Model odległości}
\label{sec:model_odleglosc}

\paragraph{Dane wejściowe.}
Z komunikatu SBS typu \texttt{MSG,3} odczytuje się szerokość $\varphi_i$
i długość $\lambda_i$ geograficzną samolotu (w stopniach) oraz jego wysokość
w stopach. Z tematu \texttt{/gps/fix} odczytuje się szerokość
$\varphi_{\mathrm{s}}$, długość $\lambda_{\mathrm{s}}$ i wysokość
$h_{\mathrm{s}}$ stacji (w metrach). Indeks $i$ oznacza samolot, a indeks
$\mathrm{s}$ -- stację.

\paragraph{Przeliczenie jednostek.}
W lotnictwie wysokość podaje się w stopach, a prędkość w węzłach. Na
jednostki układu SI przelicza się je wzorami
\begin{equation}
  h_i = \num{0.3048}\, h_{\mathrm{ft}}, \qquad
  v = \num{0.514444}\, v_{\mathrm{kt}},
  \label{eq:jednostki}
\end{equation}
\begin{objasnienia}
  \item[$h_{\mathrm{ft}}$] -- wysokość samolotu w stopach,
  \item[$h_i$] -- wysokość samolotu w metrach,
  \item[$v_{\mathrm{kt}}$] -- prędkość w węzłach,
  \item[$v$] -- prędkość w metrach na sekundę.
\end{objasnienia}

\paragraph{Odległość po powierzchni Ziemi.}
Ziemię przybliżono kulą o promieniu $R_{\mathrm{Z}} = \SI{6371}{km}$.
Odległość między stacją a punktem na powierzchni Ziemi leżącym dokładnie pod
samolotem oblicza się wzorem haversine:
\begin{equation}
  s = 2 R_{\mathrm{Z}} \arcsin \sqrt{
        \sin^2\frac{\Delta\varphi}{2}
        + \cos\varphi_{\mathrm{s}} \cos\varphi_i
          \sin^2\frac{\Delta\lambda}{2}},
  \label{eq:haversine}
\end{equation}
\begin{objasnienia}
  \item[$s$] -- odległość mierzona po powierzchni Ziemi w metrach,
  \item[$\Delta\varphi$] -- różnica szerokości geograficznych, $\Delta\varphi = \varphi_i - \varphi_{\mathrm{s}}$, w radianach,
  \item[$\Delta\lambda$] -- różnica długości geograficznych, $\Delta\lambda = \lambda_i - \lambda_{\mathrm{s}}$, w radianach,
  \item[$R_{\mathrm{Z}}$] -- średni promień Ziemi w metrach.
\end{objasnienia}
Kąty podawane w stopniach należy przed użyciem we wzorze
\eqref{eq:haversine} zamienić na radiany, mnożąc je przez $\pi/180$.

\paragraph{Odległość skośna.}
Samolot znajduje się na wysokości $h_i$, a stacja na wysokości
$h_{\mathrm{s}}$. Odległość w linii prostej od stacji do samolotu
(odległość skośną) przybliża się, traktując odległość $s$ i różnicę wysokości
jako przyprostokątne trójkąta prostokątnego:
\begin{equation}
  \rho = \sqrt{s^2 + \Delta h^2}, \qquad \Delta h = h_i - h_{\mathrm{s}},
  \label{eq:odleglosc_skosna}
\end{equation}
gdzie $\rho$ jest odległością skośną, a $\Delta h$ -- różnicą wysokości
samolotu i stacji, obie w metrach. Wzory \eqref{eq:jednostki}%
--\eqref{eq:odleglosc_skosna} są zaimplementowane w węźle wyświetlacza
stacji, który na ich podstawie pokazuje odległość do każdego samolotu.

\subsection{Uproszczenia modelu}
\label{sec:model_uproszczenia}

W modelu przyjęto następujące uproszczenia:
\begin{itemize}[nosep]
  \item Ziemia jest kulą, a nie elipsoidą -- powoduje to błąd odległości
        rzędu kilku dziesiątych procenta,
  \item we wzorze \eqref{eq:odleglosc_skosna} pominięto krzywiznę Ziemi;
        błąd rośnie wraz z odległością samolotu od stacji,
  \item wysokość barometryczna z ADS-B jest traktowana jak wysokość nad
        poziomem morza; przy ciśnieniu innym niż standardowe
        (\SI{1013.25}{hPa}) różnica może sięgać kilkudziesięciu metrów,
  \item przesunięcie zegara $\delta$ jest stałe w trakcie sesji (pominięto
        powolny dryf zegara),
  \item stacja jest nieruchoma w trakcie sesji,
  \item pominięto opóźnienie między odebraniem komunikatu ADS-B a nadaniem mu
        znacznika czasu (pojedyncze milisekundy).
\end{itemize}
Uproszczenia te są wystarczające do podglądu na żywo i do analizy zasięgu
stacji. Przy dalszym rozwoju systemu model odległości można zastąpić
dokładniejszym przeliczeniem współrzędnych do lokalnego układu kartezjańskiego
związanego ze stacją, które uwzględnia kształt Ziemi i pozwala dodatkowo
wyznaczyć kierunek obserwacji samolotu.

\subsection{Przykład obliczeniowy}
\label{sec:model_przyklad}

Przyjęto stację w punkcie $\varphi_{\mathrm{s}} = \SI{50.0}{\degree}$,
$\lambda_{\mathrm{s}} = \SI{20.0}{\degree}$, $h_{\mathrm{s}} = \SI{250}{m}$
oraz samolot w punkcie $\varphi_i = \SI{50.3169}{\degree}$,
$\lambda_i = \SI{20.49795}{\degree}$, lecący na wysokości
$h_{\mathrm{ft}} = \num{32800}$~stóp.

Z wzoru \eqref{eq:jednostki} wysokość samolotu wynosi
\begin{equation}
  h_i = \num{0.3048} \cdot \num{32800} \approx \SI{9997.4}{m},
  \qquad
  \Delta h = \num{9997.4} - 250 = \SI{9747.4}{m}.
\end{equation}
Różnice współrzędnych wynoszą $\Delta\varphi = \SI{0.3169}{\degree}$
i $\Delta\lambda = \SI{0.49795}{\degree}$. Wyrażenie pod pierwiastkiem we
wzorze \eqref{eq:haversine} ma wartość ok.\ \num{1.540e-5}, skąd
\begin{equation}
  s \approx \SI{50.00}{km}, \qquad
  \rho = \sqrt{\num{50000}^2 + \num{9747.4}^2} \approx \SI{50.94}{km}.
\end{equation}

Wpływ błędu czasu pokazuje następujące oszacowanie: samolot lecący
z prędkością $v = \SI{230}{m/s}$ w czasie $\Delta\tau$ pokonuje drogę
\begin{equation}
  \Delta x = v \, \Delta\tau .
  \label{eq:blad_czasu}
\end{equation}
Jeżeli zegar stacji myli się o \SI{10}{s} i nie zostanie poprawiony wzorem
\eqref{eq:korekta_czasu}, pozycja przypisana danej chwili będzie błędna
o ok.\ \SI{2.3}{km}. Po korekcji błąd czasu spada do ok.\ \SI{0.1}{s},
a błąd pozycji -- do ok.\ \SI{23}{m}. Uzasadnia to rejestrowanie czasu GPS
razem z danymi ADS-B.

\subsection{Weryfikacja modelu}
\label{sec:model_weryfikacja}

Poprawność modelu odległości można sprawdzić na prostych przypadkach:
dla samolotu położonego dokładnie nad stacją powinno zachodzić $s = 0$
i $\rho = \Delta h$, a dla punktów odległych o \SI{1}{\degree} szerokości
geograficznej $s \approx \SI{111.2}{km}$. Poprawność modelu czasu można
sprawdzić, porównując $\hat{\delta}$ wyznaczone w sesji z dostępem do
internetu (gdy zegar jest synchronizowany przez NTP i $\hat{\delta}$ powinno
być bliskie zera) z sesją bez dostępu do sieci. Szczegółową ocenę na
podstawie zarejestrowanych danych przedstawiono w~rozdziale~\ref{chap:badania}.
%  Wymaga w preambule: amsmath, siunitx, enumitem
% ============================================================

% Objaśnienia symboli pod wzorem ("gdzie:"): symbol w kolumnie od
% wcięcia akapitu, myślniki zawsze w jednej pionowej linii
\newenvironment{objasnienia}{%
  \noindent gdzie:
  \begin{description}[font=\normalfont, labelindent=\parindent,
    labelwidth=1.2cm, labelsep=0.2cm, leftmargin=!, align=left,
    itemsep=0pt, topsep=2pt]}{\end{description}}

\section{Model wyznaczania danych referencyjnych}
\label{sec:model}

Zarejestrowane przez stację dane mają posłużyć jako dane referencyjne, czyli
mają pozwolić stwierdzić, gdzie znajdował się dany statek powietrzny
w określonej chwili czasu. % referencja do podrozdziału.
W tym celu potrzebny jest model, który łączy trzy źródła informacji:
komunikaty ADS-B odbierane drogą radiową, pozycję stacji z odbiornika GPS oraz
zegar systemowy komputera, którym oznaczane są wszystkie zapisane wiadomości.
W niniejszym podrozdziale model ten przedstawiono najpierw w postaci słownej
(model opisowy), a następnie w postaci wzorów (model matematyczny). Na końcu
omówiono przyjęte uproszczenia i podano przykład obliczeniowy.

\subsection{Model opisowy}
\label{sec:model_opisowy}

Statek powietrzny wyposażony w transponder ADS-B kilka razy na sekundę
rozgłasza na częstotliwości \SI{1090}{MHz} komunikaty zawierające m.in.\ jego
unikalny adres ICAO, pozycję (szerokość i długość geograficzną) wyznaczoną
przez pokładowy odbiornik GNSS, wysokość barometryczną oraz prędkość
względem ziemi~\cite{sun2021}. Stacja odbiera część tych komunikatów --
niektóre giną z powodu zakłóceń, zbyt słabego sygnału lub przesłonięcia
anteny. Odebrane komunikaty są zapisywane w postaci tekstowej (format SBS)
razem ze znacznikiem czasu nadawanym przez komputer w chwili odbioru.

Komunikat ADS-B mówi, gdzie samolot jest, ale nie mówi, gdzie jest
względem stacji. Do tego potrzebna jest pozycja samej stacji, którą
podaje odbiornik GPS w komunikacie NMEA \texttt{GGA}~\cite{nmea0183}.
Znając obie pozycje, można obliczyć odległość od stacji do samolotu.

Drugim problemem jest czas. Komputer Raspberry Pi nie ma podtrzymywanego
bateryjnie zegara, a w terenie nie ma dostępu do internetu, więc jego zegar
może się mylić o nieznaną wartość. Ponieważ wszystkie zapisane wiadomości są
oznaczone właśnie tym zegarem, błąd zegara przenosi się na wszystkie dane.
Odbiornik GPS co sekundę przesyła jednak dokładny czas UTC (komunikat
\texttt{RMC}). Komputer zapisuje go razem z własnym wskazaniem zegara, dzięki
czemu po zakończeniu pomiarów można zmierzyć, o ile zegar się mylił,
i poprawić czas wszystkich wiadomości.

Granice modelowanego układu wyznacza sama stacja (antena, odbiornik RTL-SDR,
odbiornik GPS i komputer). Samoloty i satelity GPS należą do otoczenia,
a punktami styku są: łącze radiowe \SI{1090}{MHz}, łącze szeregowe
odbiornika GPS oraz plik zapisu rosbag2. Model składa się z dwóch części:
\begin{enumerate}[nosep]
  \item modelu czasu -- poprawienie znaczników czasu zegara systemowego
        na czas UTC (podrozdział~\ref{sec:model_czasu}),
  \item modelu odległości -- obliczenie odległości od stacji do samolotu
        (podrozdział~\ref{sec:model_odleglosc}).
\end{enumerate}

\subsection{Model czasu}
\label{sec:model_czasu}

Niech $t$ oznacza wskazanie zegara systemowego komputera, a $\tau$ --
odpowiadający mu rzeczywisty czas UTC. Zależność między nimi zapisano jako
\begin{equation}
  \tau = t + \delta,
  \label{eq:model_zegara}
\end{equation}
gdzie $\delta$ jest przesunięciem (błędem) zegara systemowego w sekundach.
Przyjęto, że w trakcie jednej sesji pomiarowej przesunięcie $\delta$ jest
stałe.

Węzeł obsługi GPS publikuje w temacie \texttt{/gps/time\_ref} pary
$(t_k, \tau_k)$, $k = 1, 2, \ldots, K$, gdzie $\tau_k$ jest czasem UTC
odczytanym z komunikatu \texttt{RMC}, a $t_k$ -- wskazaniem zegara
systemowego w chwili odczytu tego komunikatu. Każda para daje jeden pomiar
przesunięcia:
\begin{equation}
  \delta_k = \tau_k - t_k, \qquad k = 1, 2, \ldots, K.
  \label{eq:offset_pomiar}
\end{equation}
Pojedyncze pomiary różnią się nieco od siebie, dlatego jako wynik przyjmuje
się ich średnią arytmetyczną, a rozrzut pomiarów opisuje się odchyleniem
standardowym:
\begin{equation}
  \hat{\delta} = \frac{1}{K} \sum_{k=1}^{K} \delta_k,
  \qquad
  \sigma_{\delta} = \sqrt{\frac{1}{K-1} \sum_{k=1}^{K}
                    \bigl(\delta_k - \hat{\delta}\bigr)^2},
  \label{eq:offset_srednia}
\end{equation}
\begin{objasnienia}
  \item[$\hat{\delta}$] -- wyznaczone przesunięcie zegara w sekundach,
  \item[$\sigma_{\delta}$] -- odchylenie standardowe pomiarów przesunięcia w sekundach,
  \item[$K$] -- liczba pomiarów (przy jednym pomiarze na sekundę sesja trwająca \SI{10}{min} daje ok.~600 pomiarów).
\end{objasnienia}
Poprawiony czas UTC dowolnej zapisanej wiadomości o znaczniku $t$ wynosi
\begin{equation}
  \hat{\tau} = t + \hat{\delta}.
  \label{eq:korekta_czasu}
\end{equation}

Pomiar \eqref{eq:offset_pomiar} nie jest idealny, ponieważ znacznik $t_k$
jest nadawany dopiero po przesłaniu komunikatu łączem szeregowym i odczytaniu
go przez program. Przy prędkości \num{9600}~bit/s przesłanie jednego znaku
(10~bitów) trwa ok.\ \SI{1.04}{ms}, a komunikat NMEA ma do 82 znaków, więc
sama transmisja zajmuje do ok.\ \SI{85}{ms}. Program odczytuje port co
\SI{100}{ms}, co dodaje średnio kolejne \SI{50}{ms}. Wyznaczone przesunięcie
$\hat{\delta}$ jest więc obarczone błędem rzędu \SI{0.1}{s}, co -- jak
pokazano w podrozdziale~\ref{sec:model_przyklad} -- jest wartością
akceptowalną.

\subsection{Model odległości}
\label{sec:model_odleglosc}

\paragraph{Dane wejściowe.}
Z komunikatu SBS typu \texttt{MSG,3} odczytuje się szerokość $\varphi_i$
i długość $\lambda_i$ geograficzną samolotu (w stopniach) oraz jego wysokość
w stopach. Z tematu \texttt{/gps/fix} odczytuje się szerokość
$\varphi_{\mathrm{s}}$, długość $\lambda_{\mathrm{s}}$ i wysokość
$h_{\mathrm{s}}$ stacji (w metrach). Indeks $i$ oznacza samolot, a indeks
$\mathrm{s}$ -- stację.

\paragraph{Przeliczenie jednostek.}
W lotnictwie wysokość podaje się w stopach, a prędkość w węzłach. Na
jednostki układu SI przelicza się je wzorami
\begin{equation}
  h_i = \num{0.3048}\, h_{\mathrm{ft}}, \qquad
  v = \num{0.514444}\, v_{\mathrm{kt}},
  \label{eq:jednostki}
\end{equation}
\begin{objasnienia}
  \item[$h_{\mathrm{ft}}$] -- wysokość samolotu w stopach,
  \item[$h_i$] -- wysokość samolotu w metrach,
  \item[$v_{\mathrm{kt}}$] -- prędkość w węzłach,
  \item[$v$] -- prędkość w metrach na sekundę.
\end{objasnienia}

\paragraph{Odległość po powierzchni Ziemi.}
Ziemię przybliżono kulą o promieniu $R_{\mathrm{Z}} = \SI{6371}{km}$.
Odległość między stacją a punktem na powierzchni Ziemi leżącym dokładnie pod
samolotem oblicza się wzorem haversine:
\begin{equation}
  s = 2 R_{\mathrm{Z}} \arcsin \sqrt{
        \sin^2\frac{\Delta\varphi}{2}
        + \cos\varphi_{\mathrm{s}} \cos\varphi_i
          \sin^2\frac{\Delta\lambda}{2}},
  \label{eq:haversine}
\end{equation}
\begin{objasnienia}
  \item[$s$] -- odległość mierzona po powierzchni Ziemi w metrach,
  \item[$\Delta\varphi$] -- różnica szerokości geograficznych, $\Delta\varphi = \varphi_i - \varphi_{\mathrm{s}}$, w radianach,
  \item[$\Delta\lambda$] -- różnica długości geograficznych, $\Delta\lambda = \lambda_i - \lambda_{\mathrm{s}}$, w radianach,
  \item[$R_{\mathrm{Z}}$] -- średni promień Ziemi w metrach.
\end{objasnienia}
Kąty podawane w stopniach należy przed użyciem we wzorze
\eqref{eq:haversine} zamienić na radiany, mnożąc je przez $\pi/180$.

\paragraph{Odległość skośna.}
Samolot znajduje się na wysokości $h_i$, a stacja na wysokości
$h_{\mathrm{s}}$. Odległość w linii prostej od stacji do samolotu
(odległość skośną) przybliża się, traktując odległość $s$ i różnicę wysokości
jako przyprostokątne trójkąta prostokątnego:
\begin{equation}
  \rho = \sqrt{s^2 + \Delta h^2}, \qquad \Delta h = h_i - h_{\mathrm{s}},
  \label{eq:odleglosc_skosna}
\end{equation}
gdzie $\rho$ jest odległością skośną, a $\Delta h$ -- różnicą wysokości
samolotu i stacji, obie w metrach. Wzory \eqref{eq:jednostki}%
--\eqref{eq:odleglosc_skosna} są zaimplementowane w węźle wyświetlacza
stacji, który na ich podstawie pokazuje odległość do każdego samolotu.

\subsection{Uproszczenia modelu}
\label{sec:model_uproszczenia}

W modelu przyjęto następujące uproszczenia:
\begin{itemize}[nosep]
  \item Ziemia jest kulą, a nie elipsoidą -- powoduje to błąd odległości
        rzędu kilku dziesiątych procenta,
  \item we wzorze \eqref{eq:odleglosc_skosna} pominięto krzywiznę Ziemi;
        błąd rośnie wraz z odległością samolotu od stacji,
  \item wysokość barometryczna z ADS-B jest traktowana jak wysokość nad
        poziomem morza; przy ciśnieniu innym niż standardowe
        (\SI{1013.25}{hPa}) różnica może sięgać kilkudziesięciu metrów,
  \item przesunięcie zegara $\delta$ jest stałe w trakcie sesji (pominięto
        powolny dryf zegara),
  \item stacja jest nieruchoma w trakcie sesji,
  \item pominięto opóźnienie między odebraniem komunikatu ADS-B a nadaniem mu
        znacznika czasu (pojedyncze milisekundy).
\end{itemize}
Uproszczenia te są wystarczające do podglądu na żywo i do analizy zasięgu
stacji. Przy dalszym rozwoju systemu model odległości można zastąpić
dokładniejszym przeliczeniem współrzędnych do lokalnego układu kartezjańskiego
związanego ze stacją, które uwzględnia kształt Ziemi i pozwala dodatkowo
wyznaczyć kierunek obserwacji samolotu.

\subsection{Przykład obliczeniowy}
\label{sec:model_przyklad}

Przyjęto stację w punkcie $\varphi_{\mathrm{s}} = \SI{50.0}{\degree}$,
$\lambda_{\mathrm{s}} = \SI{20.0}{\degree}$, $h_{\mathrm{s}} = \SI{250}{m}$
oraz samolot w punkcie $\varphi_i = \SI{50.3169}{\degree}$,
$\lambda_i = \SI{20.49795}{\degree}$, lecący na wysokości
$h_{\mathrm{ft}} = \num{32800}$~stóp.

Z wzoru \eqref{eq:jednostki} wysokość samolotu wynosi
\begin{equation}
  h_i = \num{0.3048} \cdot \num{32800} \approx \SI{9997.4}{m},
  \qquad
  \Delta h = \num{9997.4} - 250 = \SI{9747.4}{m}.
\end{equation}
Różnice współrzędnych wynoszą $\Delta\varphi = \SI{0.3169}{\degree}$
i $\Delta\lambda = \SI{0.49795}{\degree}$. Wyrażenie pod pierwiastkiem we
wzorze \eqref{eq:haversine} ma wartość ok.\ \num{1.540e-5}, skąd
\begin{equation}
  s \approx \SI{50.00}{km}, \qquad
  \rho = \sqrt{\num{50000}^2 + \num{9747.4}^2} \approx \SI{50.94}{km}.
\end{equation}

Wpływ błędu czasu pokazuje następujące oszacowanie: samolot lecący
z prędkością $v = \SI{230}{m/s}$ w czasie $\Delta\tau$ pokonuje drogę
\begin{equation}
  \Delta x = v \, \Delta\tau .
  \label{eq:blad_czasu}
\end{equation}
Jeżeli zegar stacji myli się o \SI{10}{s} i nie zostanie poprawiony wzorem
\eqref{eq:korekta_czasu}, pozycja przypisana danej chwili będzie błędna
o ok.\ \SI{2.3}{km}. Po korekcji błąd czasu spada do ok.\ \SI{0.1}{s},
a błąd pozycji -- do ok.\ \SI{23}{m}. Uzasadnia to rejestrowanie czasu GPS
razem z danymi ADS-B.

\subsection{Weryfikacja modelu}
\label{sec:model_weryfikacja}

Poprawność modelu odległości można sprawdzić na prostych przypadkach:
dla samolotu położonego dokładnie nad stacją powinno zachodzić $s = 0$
i $\rho = \Delta h$, a dla punktów odległych o \SI{1}{\degree} szerokości
geograficznej $s \approx \SI{111.2}{km}$. Poprawność modelu czasu można
sprawdzić, porównując $\hat{\delta}$ wyznaczone w sesji z dostępem do
internetu (gdy zegar jest synchronizowany przez NTP i $\hat{\delta}$ powinno
być bliskie zera) z sesją bez dostępu do sieci. Szczegółową ocenę na
podstawie zarejestrowanych danych przedstawiono w~rozdziale~\ref{chap:badania}.